\gddchapter{Reflection analysis}

The role of this chapter is to provide reflections post interview and to refine the process for the next interview.

\section{Defensive self-probing}

\begin{QandA}
   \item Wait time \& silence: Did I provide the participant with time and space to voice out their opinions or did I fill the silence gap due to my own discomfort?
        \begin{answered}
           I gave the participant the necessary silence.
        \end{answered}

   \item Lexical mirroring: Identify 1-2 times where I used participants own words to deepen his narrative and answer.

        \begin{answered}
        The participant said "In the version I played..." and the researched followed up with "In the version you played...".
        \end{answered}

    \item The "Vessel trap": Did I treat the participant as a way to prove my internal bias that this system might be more immersive than traditional input?

        \begin{answered}
            In my opinion I didn't trap the participant to be the vessel for this thesis.
        \end{answered}
\end{QandA}

% \section{Sticky moments analysis}
% Analysis of the meaning of the sticky moments noted in the chapter 3


% \begin{enumerate}
%     \item \textbf{Long descriptions}
    
%     Just like in last two tests the participant had a immediate reaction when the background image for the Fishing Store location was missing.


%     \item \textbf{Not very obvious how to use items}
    
%     In the keyboard version of the game the UX of item usage was confuisng for the player. When using the items the player is suppoused to press the return key on the highlited item in the inventory. That flow seemed to be intuitive for the researcher, but the user completely missed it and it had to be explained.







%     \item \textbf{Stair naming}
    
%     This participant was visibly confused about the naming of stairs going up and down just like the rest of participants.


%     \item \textbf{Recording UX failure}
    
%     As observed before the flow of recording the voice commands on holding the button R down was not obvious to the participant.



    

%     \item \textbf{Lounge voice processing error}
    
%    Just like in test \#1 and \#3 the participant couldn't get into lounge via voice twice. The game just refused to understand that word, pointing to the fact that some words are tougher to pronounce for non-native speakers. Just as the last time she had to be manually assisted by the researcher in advancing forward.

%     \item \textbf{Pharmacy voice processing error}
    
%     In a similar fashion to the failure described above, the player couldn't enter the Pharmacy on her first attempt. Thankfully the system let her through on the second attempt.

% \end{enumerate}



\section{Refinement and iterative tweaks}

\begin{QandA}
    \item Do any questions need changes or refinements based on my observations?
    \begin{answered}
        The questions themselves seem to be working as intended. No changes are needed.
    \end{answered}

    \item Does the testing procedure need any changes based on the observations from this session?
    \begin{answered}
        I'm happy with the revered order of baseline and novelty versions of the game. I'll be keeping it going forward as it helps to reduce the modality carry-over effect that I've noted on in last 4 participants.
    \end{answered}
\end{QandA}